\documentclass[12pt]{article}
\usepackage[utf8]{inputenc}
\usepackage[T1]{fontenc}
\usepackage{hyperref}
\title{Xwela Gbe Language Resources for K-12 Curriculum Development}
\author{Compiled by Claude Sonnet 5.5 (claude-sonnet-5-5) for Andr\'as Kornai}
\date{}
\begin{document}
\maketitle

\section{Introduction}
Xwela Gbe (ISO 639-3 xwe; also called Phera, Houeda, Peda) is a Gbe language of southern Benin, spoken mainly around Lake Ahémé and nearby coastal communes. Speaker estimates are about 102,000 (Global Recordings Network) and about 106,000 (Joshua Project). It should not be confused with Western Xwla (Phla, xwl), a separate Gbe language of Benin and Togo. Few resources were found, and none for schools.

\section{Primary Educational Resources}
No primers, textbooks or school programmes for Xwela Gbe were found in the searches.

\section{Linguistic and Reference Materials}
\begin{itemize}
\item Joshua Project profile (population, scripture status): \url{https://joshuaproject.net/people_groups/15926/BN}
\item Numerals page for Western Xwla (xwl), a different language, listed here only to prevent confusion: \url{https://lingweb.eva.mpg.de/channumerals/Gbe-Western-Xwla.htm}
\item Wikipedia on the related Phla language: \url{https://en.wikipedia.org/wiki/Phla_language}
\end{itemize}
No grammar or dictionary of Xwela Gbe was found.

\section{Religious and Community Resources}
Joshua Project reports that scripture exists only as portions. Scripture Earth lists materials under the code xwe: \url{https://www.scriptureearth.org/00eng.php?iso=xwe}

\section{Digital and Audio Resources}
Global Recordings Network has recordings for xwe: \url{https://globalrecordings.net/en/language/xwe}

\section{Curriculum Development Context}
Xwela Gbe is a vigorous first language, but no school use is known, and the documentation found is thin. A curriculum project would have to start with community consultation, an agreed orthography, and fieldwork. Scripture Earth and Global Recordings Network materials could provide initial text and audio. Closely related Gbe languages with larger literatures could serve as comparison resources.
\end{document}